\chapter{Réalisation et implémentation des fonctionnalités}
\label{ch:realisation}

\section{Introduction}

Ce chapitre présente les interfaces réalisées de \nomapplication. Chaque figure correspond à un écran de l'application ; les captures sont à déposer dans \texttt{figures/captures/} (voir README du dossier).

\section{Espace public et connexion}

\subsection{Page d'accueil (landing)}

\begin{figure}[H]
    \centering
    \IfFileExists{figures/captures/landing.png}{%
        
\includegraphics[width=\textwidth]{landing.png}%
    }{%
        \fbox{\parbox{0.9\textwidth}{\centering\vspace{2cm}
        \textit{[Capture : \texttt{figures/captures/landing.png}]}
        \vspace{2cm}}}%
    }
    \caption{Page d'accueil publique \nomapplication}
    \label{fig:landing}
\end{figure}

\subsection{Page de connexion}

\begin{figure}[H]
    \centering
    \IfFileExists{figures/captures/login.png}{%
        \includegraphics[width=0.75\textwidth]{login.png}%
    }{%
        \fbox{\parbox{0.75\textwidth}{\centering\vspace{2cm}
        \textit{[Capture : \texttt{login.png}]}
        \vspace{2cm}}}%
    }
    \caption{Connexion personnel et client (CIN, n° client ou username)}
    \label{fig:login}
\end{figure}

\section{Espace agence — personnel}

\subsection{Tableau de bord}

\begin{figure}[H]
    \centering
    \IfFileExists{figures/captures/dashboard-agent.png}{%
        \includegraphics[width=\textwidth]{dashboard-agent.png}%
    }{%
        \fbox{\parbox{0.9\textwidth}{\centering\vspace{2cm}
        \textit{[Capture chef : \texttt{dashboard-agent.png}]}
        \vspace{2cm}}}%
    }
    \caption{Dashboard avec indicateurs et alertes chéquier}
    \label{fig:dashboard}
\end{figure}

\subsection{Gestion des clients}

\begin{figure}[H]
    \centering
    \IfFileExists{figures/captures/clients-liste.png}{%
        \includegraphics[width=\textwidth]{clients-liste.png}%
    }{%
        \fbox{\parbox{0.9\textwidth}{\centering\vspace{2cm}
        \textit{[Capture : \texttt{clients-liste.png}]}
        \vspace{2cm}}}%
    }
    \caption{Liste et recherche des clients}
    \label{fig:clients}
\end{figure}

\subsection{Comptes bancaires}

\begin{figure}[H]
    \centering
    \IfFileExists{figures/captures/compte-detail.png}{%
        \includegraphics[width=\textwidth]{compte-detail.png}%
    }{%
        \fbox{\parbox{0.9\textwidth}{\centering\vspace{2cm}
        \textit{[Capture : \texttt{compte-detail.png}]}
        \vspace{2cm}}}%
    }
    \caption{Fiche détail compte — solde, statut, actions}
    \label{fig:compte}
\end{figure}

\subsection{Opérations financières}

\begin{figure}[H]
    \centering
    \IfFileExists{figures/captures/operations-depot.png}{%
        \includegraphics[width=0.85\textwidth]{operations-depot.png}%
    }{%
        \fbox{\parbox{0.85\textwidth}{\centering\vspace{2cm}
        \textit{[Capture : \texttt{operations-depot.png}]}
        \vspace{2cm}}}%
    }
    \caption{Formulaire de dépôt}
    \label{fig:depot}
\end{figure}

\begin{figure}[H]
    \centering
    \IfFileExists{figures/captures/paiement-facture.png}{%
        \includegraphics[width=0.85\textwidth]{paiement-facture.png}%
    }{%
        \fbox{\parbox{0.85\textwidth}{\centering\vspace{2cm}
        \textit{[Capture : \texttt{paiement-facture.png}]}
        \vspace{2cm}}}%
    }
    \caption{Paiement de facture (LYDEC, ONEE, IAM…)}
    \label{fig:facture}
\end{figure}

\subsection{Historique et reçu PDF}

\begin{figure}[H]
    \centering
    \IfFileExists{figures/captures/recu-pdf.png}{%
        \includegraphics[width=\textwidth]{recu-pdf.png}%
    }{%
        \fbox{\parbox{0.9\textwidth}{\centering\vspace{2cm}
        \textit{[Capture : \texttt{recu-pdf.png}]}
        \vspace{2cm}}}%
    }
    \caption{Reçu d'opération avec téléchargement PDF}
    \label{fig:recu}
\end{figure}

\subsection{Commandes de chéquier et notifications}

\begin{figure}[H]
    \centering
    \IfFileExists{figures/captures/chequier-liste.png}{%
        \includegraphics[width=\textwidth]{chequier-liste.png}%
    }{%
        \fbox{\parbox{0.9\textwidth}{\centering\vspace{2cm}
        \textit{[Capture : \texttt{chequier-liste.png}]}
        \vspace{2cm}}}%
    }
    \caption{Suivi des commandes de chéquier}
    \label{fig:chequier}
\end{figure}

\begin{figure}[H]
    \centering
    \IfFileExists{figures/captures/notifications.png}{%
        \includegraphics[width=\textwidth]{notifications.png}%
    }{%
        \fbox{\parbox{0.9\textwidth}{\centering\vspace{2cm}
        \textit{[Capture : \texttt{notifications.png}]}
        \vspace{2cm}}}%
    }
    \caption{Centre de notifications (chef d'agence)}
    \label{fig:notifications}
\end{figure}

\subsection{Administration}

\begin{figure}[H]
    \centering
    \IfFileExists{figures/captures/admin-users.png}{%
        \includegraphics[width=\textwidth]{admin-users.png}%
    }{%
        \fbox{\parbox{0.9\textwidth}{\centering\vspace{2cm}
        \textit{[Capture : \texttt{admin-users.png}]}
        \vspace{2cm}}}%
    }
    \caption{Gestion des utilisateurs internes}
    \label{fig:admin}
\end{figure}

\section{Espace client (portail)}

\begin{figure}[H]
    \centering
    \IfFileExists{figures/captures/portal-dashboard.png}{%
        \includegraphics[width=\textwidth]{portal-dashboard.png}%
    }{%
        \fbox{\parbox{0.9\textwidth}{\centering\vspace{2cm}
        \textit{[Capture client CD789012 : \texttt{portal-dashboard.png}]}
        \vspace{2cm}}}%
    }
    \caption{Tableau de bord espace client}
    \label{fig:portal}
\end{figure}

\begin{figure}[H]
    \centering
    \IfFileExists{figures/captures/portal-notifications.png}{%
        \includegraphics[width=\textwidth]{portal-notifications.png}%
    }{%
        \fbox{\parbox{0.9\textwidth}{\centering\vspace{2cm}
        \textit{[Capture : \texttt{portal-notifications.png}]}
        \vspace{2cm}}}%
    }
    \caption{Notifications client après opération en agence}
    \label{fig:portal-notif}
\end{figure}

\section{Conclusion}

Ce chapitre a illustré les principales interfaces de \nomapplication\ : vitrine publique, back-office agence et portail client. L'ensemble reflète une charte graphique cohérente et une séparation claire des rôles.
